\chapter{JavaScript nel browser}

Finora abbiamo studiato il \textbf{linguaggio} JavaScript. Il linguaggio però
può essere eseguito in \textbf{ambienti ospite} (\textit{host environment})
diversi:

\begin{itemize}
  \item i \term{web browser}, per i quali JavaScript è stato originariamente
        creato;
  \item i \term{web server}, per esempio tramite il runtime Node.js.
\end{itemize}

Ogni ambiente ospite fornisce oggetti e funzioni \textbf{propri}, in aggiunta a
quelli inclusi nel nucleo del linguaggio. In questo capitolo studiamo l'ambiente
del browser.

\section{L'ambiente del browser}

\begin{figure}[H]
  \centering
  \begin{tikzpicture}[font=\Lato\footnotesize,
      grp/.style={rounded corners=4pt, line width=0.9pt, align=center,
                  minimum width=3.5cm, minimum height=2.5cm},
      it/.style={font=\FiraCodeNL\scriptsize, anchor=north, align=center}]

    \node[font=\Lato\bfseries, text=primary!75!black] (w) at (0,2.35) {\code{window}};

    \node[grp, draw=neutral!55!white, fill=neutral!7!white] (core) at (-4.0,0.4) {};
    \node[font=\Lato\scriptsize\bfseries, text=neutral!30!black, anchor=north]
      at (-4.0,1.55) {Core JavaScript};
    \node[it, text=neutral!25!black] at (-4.0,1.1)
      {Object \quad Function\\ Array \quad Map \quad Set\\ \ldots};

    \node[grp, draw=primary!60!black, fill=primary!7!white] (dom) at (0,0.4) {};
    \node[font=\Lato\scriptsize\bfseries, text=primary!70!black, anchor=north]
      at (0,1.55) {DOM};
    \node[font=\Lato\scriptsize, text=primary!55!black, anchor=north]
      at (0,1.22) {Document Object Model};
    \node[it, text=primary!60!black] at (0,0.85) {document};

    \node[grp, draw=accent!60!black, fill=accent!8!white] (bom) at (4.0,0.4) {};
    \node[font=\Lato\scriptsize\bfseries, text=accent!70!black, anchor=north]
      at (4.0,1.55) {BOM};
    \node[font=\Lato\scriptsize, text=accent!55!black, anchor=north]
      at (4.0,1.22) {Browser Object Model};
    \node[it, text=accent!60!black] at (4.0,0.85)
      {navigator \quad screen\\ location \quad history\\ frames \quad \ldots};

    \begin{scope}[on background layer]
      \node[rounded corners=5pt, draw=primary!70!black, line width=1.1pt,
            fill=primary!2!white, fit=(w)(core)(dom)(bom), inner sep=9pt] {};
    \end{scope}
  \end{tikzpicture}
  \caption{L'ambiente del browser: il nucleo del linguaggio, il DOM e il BOM.}
  \label{fig:browser-env}
\end{figure}

\subsection{L'oggetto window}

L'ambiente del browser ha un oggetto \guillemotleft radice\guillemotright\ chiamato \term{window},
che svolge due ruoli:

\begin{enumerate}
  \item è l'\textbf{oggetto globale} per il codice JavaScript;
  \item rappresenta la \textbf{finestra del browser} e fornisce i metodi per
        controllarla.
\end{enumerate}

Le funzioni dichiarate a livello globale (e le variabili dichiarate con
\code{var}) diventano proprietà dell'oggetto \code{window}.

\begin{affianco}[0.55][c]
\begin{lstlisting}[style=js, name=window.js]
let msg = "Hello!";

function greet(name) {
  console.log(`${msg}, ${name}!`);
}

greet("Web Technologies");       // Hello, Web Technologies!
window.greet("window object");   // Hello, window object!
\end{lstlisting}
\risultato
\begin{immagine}[\code{greet} è una proprietà di \code{window};\\
                 \code{msg}, dichiarata con \code{let}, no.]
\begin{tikzpicture}
  \mchiave{1.4cm}
  \oggetto{w}{(0,0)}{3.0cm}{window}
    {greet/function, document/\ldots, location/\ldots, \ldots/}
  \variabile{m}{(0.45,-3.3)}{msg}{"Hello!"}
  \scopo{1}{g}{let e const globali}{(m)(m-nome)}
  \node[mnota, anchor=west, text=merr] at ([xshift=4pt]g.east)
    {\code{window.msg}\\ vale \code{undefined}};
\end{tikzpicture}
\end{immagine}
\end{affianco}

Le variabili dichiarate con \code{let} e \code{const} sono anch'esse globali,
ma vivono in uno scope separato e \textbf{non} diventano proprietà di
\code{window}: \code{window.msg} vale \code{undefined}.

\subsection{Il DOM e il BOM}

Il \term{DOM} (\textit{Document Object Model}) rappresenta il contenuto del
documento corrente. L'oggetto \code{document} è il punto d'ingresso principale
alla pagina e fornisce i modi per accedere ai contenuti e modificarli.

\begin{affianco}[0.6][c]
\begin{lstlisting}[style=html, name=dom-primo.html]
<body>
  <p>Hello!</p>
  <script>
    document.body.style.background = "pink";
    document.querySelector("p").innerText = "Hello DOM!";
  </script>
</body>
\end{lstlisting}
\risultato
\begin{browser}[dom-primo.html]
  \vspace*{-6pt}\noindent\makebox[\linewidth][l]{\hspace*{-8pt}%
  \colorbox{css-pink}{\makebox[\dimexpr\linewidth+16pt-2\fboxsep][l]{%
    \rule[-14pt]{0pt}{36pt}Hello DOM!}}}\vspace*{-6pt}
\end{browser}
\wtnota{Sfondo rosa e testo cambiato: il DOM\\ è stato modificato dallo script.}
\end{affianco}

Il \term{BOM} (\textit{Browser Object Model}) comprende invece gli oggetti
aggiuntivi forniti dal browser per lavorare \textbf{con il browser stesso},
non con il contenuto del documento.

\begin{lstlisting}[style=js, name=bom.js]
// location contiene informazioni sulla URL del documento corrente
console.log(location.href);        // http://localhost:3000/07-JavaScript...

// screen contiene informazioni sul display usato dall'utente
console.log(screen.availWidth);    // 2048

// navigator contiene dettagli sul browser e sulla piattaforma
console.log(navigator.userAgent);  // Mozilla/5.0 (Windows NT 10.0...

// history rappresenta la cronologia di navigazione
history.back();                    // come premere "indietro" nel browser
\end{lstlisting}

\section{Il DOM}

Quando abbiamo studiato i selettori CSS avevamo accennato al fatto che i
documenti HTML si possono vedere come \textbf{alberi}: parlavamo di discendenti,
figli, fratelli. Il DOM \textbf{formalizza} quell'intuizione.

\begin{itemize}
  \item Ogni tag HTML è un oggetto nel DOM.
  \item I tag annidati sono \textbf{figli} di quelli che li racchiudono.
  \item Anche il \textbf{contenuto testuale} di un tag è un oggetto, ed è figlio
        del tag.
  \item Gli \textbf{attributi} di un tag sono accessibili come proprietà
        dell'oggetto corrispondente.
\end{itemize}

\begin{figure}[H]
  \centering
  \begin{minipage}[c]{0.38\linewidth}
\begin{lstlisting}[style=html, basicstyle=\linespread{1.1}\scriptsize\FiraCodeNL]
<!DOCTYPE html>
<html>
  <head>
    <meta charset="utf-8">
    <title>DOM</title>
  </head>
  <body>
    <table>
      <tbody>
        <tr>
          <th>Exam</th>
          <th>Grade</th>
        </tr>
        <tr>
          <td>
            <a href="/wt">Web T.</a>
          </td>
          <td>30</td>
        </tr>
      </tbody>
    </table>
  </body>
</html>
\end{lstlisting}
  \end{minipage}\hfill
  \begin{minipage}[c]{0.59\linewidth}
    \centering
    % Il disegno e' piu' largo della colonna: si lascia a TikZ la spaziatura
    % corretta (niente nodi sovrapposti) e si riscala il risultato.
    \resizebox{\linewidth}{!}{%
    \begin{tikzpicture}[font=\FiraCodeNL\scriptsize,
        level distance=0.95cm,
        n/.style={rounded corners=2pt, draw=primary!60!black, line width=0.7pt,
                  fill=primary!8!white, inner sep=2pt},
        t/.style={rounded corners=2pt, draw=green-gradient-6, line width=0.7pt,
                  fill=green-gradient-1!45!white, inner sep=2pt, align=center,
                  font=\FiraCodeNL\tiny},
        edge from parent/.style={draw=neutral!50!black, line width=0.55pt,
                                 edge from parent path=
                                 {(\tikzparentnode.south) -- (\tikzchildnode.north)}},
        level 1/.style={sibling distance=5.4cm},
        level 2/.style={sibling distance=2.7cm},
        level 3/.style={sibling distance=2.7cm},
        level 4/.style={sibling distance=2.7cm},
        level 5/.style={sibling distance=1.35cm}]
      \node[n] {<html>}
        child { node[n] {<head>}
          child { node[n] {<meta>} }
          child { node[n] {<title>}
            child { node[t] {\#text\\ DOM} } } }
        child { node[n] {<body>}
          child { node[n] {<table>}
            child { node[n] {<tbody>}
              child { node[n] {<tr>}
                child { node[n] {<th>} }
                child { node[n] {<th>} } }
              child { node[n] {<tr>}
                child { node[n] {<td>}
                  child { node[n] {<a>} } }
                child { node[n] {<td>} } } } } };
    \end{tikzpicture}}
  \end{minipage}
  \caption{Un documento HTML e il corrispondente albero DOM.}
  \label{fig:dom-example}
\end{figure}

\warning{Anche gli spazi sono nodi}{
  La figura è una semplificazione. Il DOM reale include anche gli \textbf{spazi}
  e gli \textbf{a-capo} fra i tag, sotto forma di nodi di testo. Per brevità
  questi nodi vengono di norma omessi nei disegni, ma esistono: è il motivo per
  cui \code{document.body.firstChild} molto spesso non è l'elemento che ci si
  aspetta, ma un nodo di testo contenente \code{"\textbackslash n\textbackslash t"}.}

Anche i \textbf{commenti} HTML, pur non essendo visualizzati, vengono analizzati
e aggiunti al DOM come speciali nodi commento.

\subsection{Tipi di nodo}

La specifica del DOM definisce 12 tipi di nodo. Nella pratica se ne usano
quattro:

\begin{center}
\renewcommand{\arraystretch}{1.35}
\begin{tabular}{@{}l@{\hspace{1.1cm}}>{\raggedright\arraybackslash}p{8.4cm}@{}}
  \textbf{Tipo} & \textbf{Rappresenta} \\
  nodo \code{document} & la radice e il punto d'ingresso del DOM \\
  nodi \textit{Element} & gli elementi HTML \\
  nodi \textit{Text}    & il contenuto testuale \\
  nodi \textit{Comment} & i commenti \\
\end{tabular}
\end{center}

\section{Trovare elementi}

L'oggetto \code{document} fornisce diversi metodi per selezionare elementi.

\begin{affianco}[0.62][c]
\codicepiccolo
\begin{lstlisting}[style=html, name=trova.html]
<ul>
  <li id="a">A</li>
  <li class="foo bar">B</li>
  <li name="c">C</li>
</ul>
<script>
  console.log(document.getElementById("a"));
  console.log(document.getElementsByClassName("foo"));  // HTMLCollection (1)
  console.log(document.getElementsByName("c"));         // NodeList (1)
  console.log(document.getElementsByTagName("li"));     // HTMLCollection (3)
</script>
\end{lstlisting}
\risultato
\begin{console}
<li id="a">A</li>
HTMLCollection [ li.foo.bar ]
NodeList [ li ]
HTMLCollection(3) [ li#a,
  li.foo.bar, li ]
\end{console}
\end{affianco}

I metodi \code{querySelector(css)} e \code{querySelectorAll(css)} sono i più
versatili, perché prendono in ingresso un \textbf{selettore CSS}: tutto ciò che
abbiamo imparato nel capitolo sui CSS torna utile qui.

\begin{itemize}
  \item \code{querySelector} restituisce solo il \textbf{primo} elemento che
        corrisponde al selettore;
  \item \code{querySelectorAll} restituisce una \textbf{lista} di
        corrispondenze;
  \item tutti i metodi restituiscono \code{null} quando non trovano nulla.
\end{itemize}

\begin{affianco}[0.62][c]
\codicepiccolo
\begin{lstlisting}[style=html, name=queryselector.html]
<ul>
  <li id="a">A</li> <li class="foo bar">B</li> <li name="c">C</li>
</ul>
<script>
  console.log(document.querySelector("ul > li"));         // primo <li>
  console.log(document.querySelectorAll("ul > li"));      // NodeList (3)
  console.log(document.querySelector("ul > li.myclass")); // null
</script>
\end{lstlisting}
\risultato
\begin{console}
<li id="a">A</li>
NodeList(3) [ li#a, li.foo.bar,
  li ]
null
\end{console}
\end{affianco}

\subsection{Collezioni vive}

Tutti i metodi della famiglia \code{getElementsBy*} restituiscono una
\code{HTMLCollection}, detta anche \term{collezione viva}: si aggiorna
automaticamente e riflette sempre lo stato corrente del documento.

\begin{esempio}{Una collezione che cambia da sola}
\begin{affianco}[0.6][c]
\begin{lstlisting}[style=html, name=live.html]
<a href="/">First link</a>
<script>
  let links = document.getElementsByTagName("a");
  console.log(links);   // HTMLCollection (1)
</script>

<a href="/">Second Link</a> <a href="/">Third Link</a>
<script>
  console.log(links);   // HTMLCollection (3), senza nuova ricerca!
</script>
\end{lstlisting}
\risultato
\begin{immagine}[La stessa collezione, in due momenti:\\
                 si aggiorna mentre il documento cresce.]
\begin{tikzpicture}
  \noindici
  \variabile[minimum width=0.55cm]{l1}{(0,0)}{links}{}
  \celle[minimum width=0.6cm]{c1}{(1.0,0)}{<a>}
  \riferimento{l1}{-- (c1-0.west)}
  \node[mnota, anchor=west] at ([xshift=6pt]c1-box.east) {dopo il 1\textsuperscript{o} script};
  \variabile[minimum width=0.55cm]{l2}{(0,-1.0)}{links}{}
  \celle[minimum width=0.6cm]{c2}{(1.0,-1.0)}{<a>, <a>, <a>}
  \riferimento{l2}{-- (c2-0.west)}
  \node[draw=accent!60!black, line width=1.1pt, fit=(c2-1)(c2-2), inner sep=0pt] {};
  \node[mnota, anchor=west] at ([xshift=6pt]c2-box.east) {dopo il 2\textsuperscript{o}};
\end{tikzpicture}
\end{immagine}
\end{affianco}

La variabile \code{links} non è stata riassegnata, eppure contiene tre elementi:
la collezione si è aggiornata da sé.
\end{esempio}

\subsection{Cercare dentro un elemento}

Quando i metodi di ricerca sono invocati su \code{document}, cercano
nell'\textbf{intero} documento HTML. Se gli elementi che ci servono sono
discendenti di un altro elemento, può essere più efficiente invocare i metodi di
ricerca \textbf{sull'elemento antenato}.

\begin{lstlisting}[style=js, name=ricerca-locale.js]
let sidebar = document.getElementById("sidebar");
let sidebarLinks = sidebar.querySelectorAll("a");   // cerca solo nella sidebar
\end{lstlisting}

\subsection{Riepilogo}

\begin{center}
\renewcommand{\arraystretch}{1.35}
\begin{tabular}{@{}l@{\hspace{0.8cm}}l@{\hspace{0.7cm}}c@{\hspace{0.7cm}}c@{}}
  \textbf{Metodo} & \textbf{Cerca per\ldots} & \textbf{Su elementi?} &
  \textbf{Collez. viva?} \\
  \code{querySelector}          & selettore CSS   & sì & no \\
  \code{querySelectorAll}       & selettore CSS   & sì & no \\
  \code{getElementById}         & attributo id    & no & no \\
  \code{getElementsByName}      & attributo name  & no & sì \\
  \code{getElementsByTagName}   & nome del tag o \code{*} & sì & sì \\
  \code{getElementsByClassName} & attributo class & sì & sì \\
\end{tabular}
\end{center}

\info{Quale usare}{
  Nella pratica \code{querySelector} e \code{querySelectorAll} coprono
  praticamente ogni esigenza e hanno il vantaggio di usare la stessa sintassi
  dei CSS. Gli altri metodi restano utili quando serve una collezione viva,
  oppure per \code{getElementById}, che è il più veloce di tutti.}

\section{Navigare il DOM}

Gli elementi HTML più esterni sono disponibili anche come proprietà
dell'oggetto \code{document}:

\begin{itemize}
  \item \code{document.documentElement} per l'elemento \tg{html};
  \item \code{document.head} per l'elemento \tg{head};
  \item \code{document.body} per l'elemento \tg{body}.
\end{itemize}

In generale, ogni nodo del DOM contiene riferimenti al proprio genitore, ai
fratelli e ai figli. Questi riferimenti sono in \textbf{sola lettura}.

\begin{figure}[H]
  \centering
  \begin{tikzpicture}[font=\FiraCodeNL\scriptsize,
      e/.style={rounded corners=3pt, draw=primary!65!black, line width=1pt,
                fill=primary!12!white, minimum width=2.2cm, minimum height=0.85cm},
      r/.style={rounded corners=2pt, draw=neutral!55!white, line width=0.7pt,
                fill=neutral!8!white, minimum width=2.4cm, minimum height=0.6cm},
      a/.style={-{Stealth[length=5pt,width=4pt]}, line width=0.8pt,
                draw=accent!70!black}]

    \node[e] (el) at (0,0) {<elem>};

    \node[r] (par) at (0,1.6)  {parentNode};
    \node[r] (ps)  at (-3.2,0) {previousSibling};
    \node[r] (ns)  at (3.2,0)  {nextSibling};
    \node[r] (fc)  at (-2.6,-1.6) {firstChild};
    \node[r] (lc)  at (2.6,-1.6)  {lastChild};
    \node[r] (cn)  at (0,-2.6)    {childNodes};

    \draw[a] (el) -- (par);
    \draw[a] (el) -- (ps);
    \draw[a] (el) -- (ns);
    \draw[a] (el) -- (fc);
    \draw[a] (el) -- (lc);
    \draw[a] (el) -- (cn);
  \end{tikzpicture}
  \caption{Le proprietà di navigazione fra i nodi del DOM.}
  \label{fig:dom-nav}
\end{figure}

\begin{esempio}{Navigazione fra nodi}
\begin{affianco}[0.6][c]
\codicepiccolo
\begin{lstlisting}[style=html, name=navigazione.html]
<body>
  <main>
    <h1>Hello DOM!</h1>
    <p>DOM is <em>FUN</em>!</p>
  </main>
  <script>
    console.log(document.body.firstChild.nextSibling); // elemento main
    console.log(document.body.childNodes[1]);          // elemento main
    console.log(document.querySelector("em").parentNode); // elemento p
  </script>
</body>
\end{lstlisting}
\risultato
\begin{immagine}[I figli di \code{body}: il primo è il nodo di testo\\
                 con l'a-capo e l'indentazione.]
\begin{tikzpicture}[
    a/.style={-{Stealth[length=4pt,width=3.5pt]}, line width=0.75pt,
      draw=accent!70!black}]
  \node[mdom] (b) at (2.2,0) {<body>};
  \node[mtesto] (t0) at (0,-1.4) {\#text};
  \node[mdom]   (m)  at (2.2,-1.4) {<main>};
  \node[mtesto] (t1) at (4.4,-1.4) {\#text};
  \foreach \n in {t0,m,t1} \draw[mramo] (b.south) -- (\n.north);
  \draw[a] (t0.east) -- node[mnota, above]{\code{nextSibling}} (m.west);
  \node[mnota, anchor=north] at ([yshift=-2pt]t0.south) {\code{firstChild}};
  \node[mnota, anchor=north] at ([yshift=-2pt]m.south) {\code{childNodes[1]}};
\end{tikzpicture}
\end{immagine}
\end{affianco}

Si noti \code{firstChild.nextSibling}: il primo figlio di \code{body} è il nodo
di testo contenente l'a-capo e l'indentazione, non \code{<main>}.
\end{esempio}

\subsection{Navigare fra soli elementi}

Spesso vogliamo \textbf{ignorare} i nodi di testo e di commento e concentrarci
sui soli elementi. Esistono proprietà di navigazione dedicate:

\begin{center}
\renewcommand{\arraystretch}{1.35}
\begin{tabular}{@{}l@{\hspace{1.6cm}}l@{}}
  \textbf{Tutti i nodi} & \textbf{Solo elementi} \\
  \code{parentNode}      & \code{parentElement} \\
  \code{previousSibling} & \code{previousElementSibling} \\
  \code{nextSibling}     & \code{nextElementSibling} \\
  \code{firstChild}      & \code{firstElementChild} \\
  \code{lastChild}       & \code{lastElementChild} \\
  \code{childNodes}      & \code{children} \\
\end{tabular}
\end{center}

\info{Usare quasi sempre la colonna di destra}{
  Nella pratica le proprietà \code{*Element*} sono quelle che si usano: evitano
  di dover saltare a mano i nodi di testo generati dall'indentazione del
  sorgente HTML, che sono un dettaglio di formattazione e non contenuto.}

\section{Proprietà dei nodi}

\subsection{Nome e tipo}

\code{nodeName} e \code{tagName} danno informazioni sul tipo di un nodo. Sono in
sola lettura.

\begin{affianco}[0.6][c]
\begin{lstlisting}[style=html, name=nodename.html]
<body>
  <h1 id="h">Hello <em>DOM!</em></h1>
  <!-- commento -->
  <script>
    let body = document.body;
    console.log(body.firstChild.nodeName);          // #text
    console.log(body.firstChild.tagName);           // undefined
    console.log(body.childNodes[3].nodeName);       // #comment
    console.log(body.firstElementChild.nodeName);   // H1
    console.log(body.firstElementChild.tagName);    // H1
  </script>
</body>
\end{lstlisting}
\risultato
\begin{immagine}[\code{childNodes} di \code{body}: gli elementi\\
                 sono intervallati da nodi di testo.]
\begin{tikzpicture}
  \foreach \t/\st [count=\i from 0] in {\#text/mtesto, H1/mdom, \#text/mtesto,
      \#comment/mcommento, \#text/mtesto, SCRIPT/mdom} {
    \node[\st, anchor=west, minimum width=2.1cm] (n\i) at (0.9,-\i*0.6) {\t};
    \node[mindice, anchor=east] at (0.8,-\i*0.6) {[\i]};
  }
  \node[mnota, anchor=west] at ([xshift=4pt]n0.east) {\code{firstChild}};
  \node[mnota, anchor=west] at ([xshift=4pt]n1.east) {\code{firstElementChild}};
  \node[mnota, anchor=west] at ([xshift=4pt]n3.east) {\code{childNodes[3]}};
\end{tikzpicture}
\end{immagine}
\end{affianco}

\subsection{Attributi}

Gli attributi \textbf{standard} vengono analizzati e resi disponibili come
proprietà degli elementi corrispondenti. Quelli non standard no, e vanno letti
con \code{getAttribute()}.

\begin{affianco}[0.6][c]
\begin{lstlisting}[style=html, name=attributi.html]
<body>
  <input id="x" class="inp ok" name="field" custom="y">
  <script>
    let i = document.body.firstElementChild;
    console.log(i.id);          // x
    console.log(i.classList);   // DOMTokenList ["inp", "ok"]
    console.log(i.name);        // field
    console.log(i.custom);      // undefined: custom non e' standard
    console.log(i.getAttribute("custom"));   // y
  </script>
</body>
\end{lstlisting}
\risultato
\begin{immagine}[Gli attributi standard diventano proprietà;\\
                 \code{custom} resta solo un attributo.]
\begin{tikzpicture}[
    a/.style={-{Stealth[length=4pt,width=3.5pt]}, line width=0.75pt,
      draw=neutral!45!black}]
  \mchiave{1.05cm}
  \oggetto[fill=neutral!12!white]{h}{(0,0)}{2.6cm}{attributi HTML}
    {id/"x", class/"inp ok", name/"field", custom/"y"}
  \oggetto{d}{(3.6,0)}{2.6cm}{proprietà DOM}
    {id/"x", classList/[inp ok], name/"field", custom/undefined}
  \foreach \i in {1,2,3} \draw[a] (h-\i.east) -- (d-\i.west);
  \node at ($(h-4.east)!0.5!(d-4.west)$) {\merr};
\end{tikzpicture}
\end{immagine}
\end{affianco}

Queste proprietà sono anche \textbf{scrivibili}:

\begin{lstlisting}[style=html, name=attributi-scrittura.html]
<body>
  <input id="x" class="inp ok" name="field" custom="y">
  <script>
    let i = document.body.firstElementChild;
    i.id = "w";
    i.classList.remove("ok");
    i.classList.add("err");
    i.setAttribute("custom", "k");
    i.removeAttribute("name");
  </script>
</body>

<!-- l'input dopo l'esecuzione dello script -->
<input id="w" class="inp err" custom="k">
\end{lstlisting}

\subsection{innerHTML e outerHTML}

\begin{itemize}
  \item \code{innerHTML} permette di accedere al codice HTML \textbf{contenuto}
        in un \code{HTMLElement};
  \item \code{outerHTML} permette di accedere al codice HTML
        \textbf{dell'intero} elemento.
\end{itemize}

\begin{affianco}[0.6][c]
\begin{lstlisting}[style=html, name=innerhtml.html]
<body>
  <h1 id="h">Hello <em>DOM!</em></h1>
  <script>
    let body = document.body;
    console.log(body.firstElementChild.innerHTML);
    // Hello <em>DOM!</em>
    console.log(body.firstChild.innerHTML);
    // undefined (un nodo di testo non e' un HTMLElement)
    console.log(body.firstElementChild.outerHTML);
    // <h1 id="h">Hello <em>DOM!</em></h1>
  </script>
</body>
\end{lstlisting}
\risultato
\begin{immagine}[\code{innerHTML} è il contenuto,\\
                 \code{outerHTML} l'elemento intero.]
\begin{tikzpicture}
  \node[mcod, anchor=base west, inner sep=1pt] (n1) at (0,0) {<h1 id="h">};
  \node[mcod, anchor=base west, inner sep=1pt, fill=accent!12!white]
    (n2) at (n1.base east) {Hello <em>DOM!</em>};
  \node[mcod, anchor=base west, inner sep=1pt] (n3) at (n2.base east) {</h1>};
  \draw[wtbrace, decoration={mirror=false}] ([yshift=2pt]n2.north west) -- ([yshift=2pt]n2.north east)
    node[midway, above=5pt, mnota] {\code{innerHTML}};
  \draw[wtbrace] ([yshift=-2pt]n1.south west) -- ([yshift=-2pt]n3.south east)
    node[midway, below=5pt, mnota] {\code{outerHTML}};
\end{tikzpicture}
\end{immagine}
\end{affianco}

La cosa interessante è che sono \textbf{scrivibili}: assegnando loro un nuovo
valore si cambia il contenuto della pagina.

\begin{affianco}[0.58][c]
\begin{lstlisting}[style=html, name=innerhtml-scrittura.html]
<body>
  <h1 id="h">Hello <em>DOM!</em></h1>
  <script>
    let elem = document.body.firstElementChild;

    elem.innerHTML = "DOM: <em>Wow!</em>";
    // cambiato il contenuto dell'elemento <h1>

    elem.outerHTML = `<p>${elem.innerHTML}</p>`;
    // ora l'<h1> e' diventato un <p>, con lo stesso contenuto
  </script>
</body>
\end{lstlisting}
\risultato
\begin{browser}[all'inizio]
  {\Lora\bfseries\Huge Hello \textit{DOM!}}
\end{browser}
\begin{browser}[dopo innerHTML]
  {\Lora\bfseries\Huge DOM: \textit{Wow!}}
\end{browser}
\begin{browser}[dopo outerHTML]
  DOM: \textit{Wow!}
\end{browser}
\wtnota{L'ultimo è un paragrafo: niente più stile da titolo.}
\end{affianco}

\error{innerHTML e input dell'utente}{
  Assegnare a \code{innerHTML} una stringa che contiene, anche solo in parte,
  dati forniti dall'utente è il modo classico di aprire una vulnerabilità di
  tipo \textit{Cross-Site Scripting} (XSS): il browser interpreta come markup
  ciò che gli si passa, quindi anche eventuali tag \tg{script} o attributi
  \code{onerror}. Quando serve inserire solo testo si usa \code{textContent} o
  \code{innerText}, che non interpretano nulla. Torneremo sull'argomento nel
  capitolo sulla sicurezza.}

\section{Modificare la struttura del DOM}

Non siamo limitati a modificare gli elementi esistenti: possiamo creare
dinamicamente nuovi elementi, rimuoverne di esistenti e cambiare la struttura
dell'albero. È la chiave per costruire pagine \guillemotleft vive\guillemotright\ e reattive.

\begin{affianco}[0.52][c]
\codicepiccolo
\begin{lstlisting}[style=html, name=crea.html]
<body>
  <script>
    let heading = document.createElement("h1");  // nuovo <h1>
    // il nuovo <h1> non e' ancora nel DOM

    let title = document.createTextNode("Wow!"); // nodo di testo
    // nemmeno il nodo di testo e' ancora nel DOM

    heading.append(title);           // il testo nell'<h1>
    document.body.prepend(heading);  // l'<h1> in cima al <body>
  </script>
</body>
\end{lstlisting}
\risultato
\begin{immagine}[L'albero DOM prima e dopo lo script.]
\begin{tikzpicture}[font=\FiraCodeNL\scriptsize,
      c/.style={rounded corners=2pt, draw=neutral!50!white, line width=0.7pt,
                fill=neutral!7!white, minimum width=1.6cm, minimum height=0.55cm},
      cn/.style={c, draw=accent!60!black, fill=accent!14!white}]

    \node[font=\Lato\scriptsize\itshape, text=neutral!35!black] at (-1.6,1.5)
      {prima};
    \node[font=\Lato\scriptsize\itshape, text=neutral!35!black] at (3.6,1.5)
      {dopo};

    % prima
    \node[c] (b1) at (-1.6,0.9) {<body>};
    \node[c] (s1) at (-1.6,-0.2) {<script>};
    \draw[draw=neutral!50!black, line width=0.55pt] (b1.south) -- (s1.north);

    % dopo
    \node[c]  (b2) at (3.6,0.9) {<body>};
    \node[cn] (h2) at (2.6,-0.2) {<h1>};
    \node[c]  (s2) at (4.6,-0.2) {<script>};
    \node[cn] (t2) at (2.6,-1.3) {\#text Wow!};
    \draw[draw=neutral!50!black, line width=0.55pt] (b2.south) -- (h2.north);
    \draw[draw=neutral!50!black, line width=0.55pt] (b2.south) -- (s2.north);
    \draw[draw=neutral!50!black, line width=0.55pt] (h2.south) -- (t2.north);
  \end{tikzpicture}
\end{immagine}
\end{affianco}

\subsection{Dove inserire i nodi}

\begin{esempio}{Costruire una lista dinamicamente}
\begin{lstlisting}[style=html, name=lista-dinamica.html]
<body>
  <ul>
    <li class="kw">CSS</li>
  </ul>
</body>
\end{lstlisting}

\begin{affianco}[0.5][c]
\begin{lstlisting}[style=js, name=lista-dinamica.js]
let ul   = document.querySelector("ul");
let html = ul.querySelector("li").cloneNode();
let js   = document.createElement("li");

html.innerHTML = "HTML";
js.innerHTML   = "JavaScript";

ul.before("Keywords:");
ul.prepend(html);
ul.append(js);
ul.after("End of keyword list");
\end{lstlisting}
\risultato
\begin{immagine}[I quattro punti di inserimento rispetto\\ a un elemento.]
\begin{tikzpicture}[font=\FiraCodeNL\footnotesize,
      li/.style={draw=primary!60!black, fill=primary!10!white, line width=0.7pt,
                 minimum width=3.6cm, minimum height=0.5cm, anchor=north},
      lab/.style={font=\FiraCodeNL\scriptsize, text=accent!70!black, anchor=west},
      a/.style={-{Stealth[length=4pt,width=3.5pt]}, line width=0.7pt,
                draw=accent!70!black}]

    \node[li] (l1) at (0,0)          {<li>1</li>};
    \node[li, below=1pt of l1] (l2)  {<li>2</li>};
    \node[li, below=1pt of l2] (l3)  {<li>3</li>};

    \begin{scope}[on background layer]
      \node[rounded corners=3pt, draw=neutral!50!black, line width=0.8pt,
            fit=(l1)(l2)(l3), inner sep=7pt] (ul) {};
    \end{scope}

    \draw[a] (2.7,0.85) -- (0.9,0.85);   \node[lab] at (2.8,0.85) {ul.before(...)};
    \draw[a] (2.7,0.15) -- (1.85,0.15);  \node[lab] at (2.8,0.15) {ul.prepend(...)};
    \draw[a] (2.7,-1.55) -- (1.85,-1.55);\node[lab] at (2.8,-1.55) {ul.append(...)};
    \draw[a] (2.7,-2.3) -- (0.9,-2.3);   \node[lab] at (2.8,-2.3) {ul.after(...)};
  \end{tikzpicture}
\end{immagine}
\end{affianco}

Il risultato, nel DOM e nel browser:

\begin{affianco}[0.5][c]
\begin{lstlisting}[style=html, name=risultato.html]
<body>
  Keywords:
  <ul>
    <li class="kw">HTML</li>
    <li class="kw">CSS</li>
    <li>JavaScript</li>
  </ul>
  End of keyword list
</body>
\end{lstlisting}
\risultato
\begin{browser}[lista-dinamica.html]
  Keywords:
  \begin{itemize}[leftmargin=18pt, itemsep=0pt, topsep=2pt]
    \item HTML
    \item CSS
    \item JavaScript
  \end{itemize}
  End of keyword list
\end{browser}
\end{affianco}
\end{esempio}

\begin{esempio}{Spostare e sostituire}
\begin{affianco}[0.55][c]
\begin{lstlisting}[style=js, name=sostituisci.js]
let ul = document.querySelector("ul");
let ol = document.createElement("ol");

ol.append(ul.firstElementChild);   // sposta il <li> nell'<ol>
ul.replaceWith(ol);                // l'<ol> sostituisce l'<ul>
\end{lstlisting}
\risultato
\begin{immagine}[Il primo \code{<li>} viene spostato nell'\code{<ol>},\\
                 che prende il posto dell'\code{<ul>}.]
\begin{tikzpicture}[
    a/.style={-{Stealth[length=5pt,width=4pt]}, line width=0.8pt,
      draw=neutral!45!black}]
  \node[mnota] at (0.6,0.6) {prima};
  \node[mdom] (u) at (0.6,0) {<ul>};
  \node[mnuovo] (l1) at (0,-1) {<li>};
  \node[mdom] (l2) at (1.2,-1) {<li>};
  \draw[mramo] (u.south) -- (l1.north); \draw[mramo] (u.south) -- (l2.north);
  \node[mnota] at (4.0,0.6) {dopo};
  \node[mnuovo] (o) at (4.0,0) {<ol>};
  \node[mnuovo] (m1) at (4.0,-1) {<li>};
  \draw[mramo] (o.south) -- (m1.north);
  \draw[a] (1.9,-0.5) -- (3.2,-0.5);
\end{tikzpicture}
\end{immagine}
\end{affianco}

Da notare: \code{append} non copia il nodo, lo \textbf{sposta}. Un nodo può
trovarsi in un solo punto dell'albero.
\end{esempio}

\subsection{Riepilogo dei metodi}

\begin{center}
\renewcommand{\arraystretch}{1.35}
\begin{tabular}{@{}l@{\hspace{0.9cm}}>{\raggedright\arraybackslash}p{7.4cm}@{}}
  \textbf{Metodo} & \textbf{Descrizione} \\
  \code{document.createElement(tag)}   & crea un elemento con il tag dato \\
  \code{document.createTextNode(val)}  & crea un nodo di testo (usato di rado) \\
  \code{elem.cloneNode(deep)}          & clona l'elemento; con \code{deep=true}
                                         anche i discendenti \\
  \code{node.append(...)}              & inserisce dentro \code{node}, in fondo \\
  \code{node.prepend(...)}             & inserisce dentro \code{node}, all'inizio \\
  \code{node.before(...)}              & inserisce subito prima di \code{node} \\
  \code{node.after(...)}               & inserisce subito dopo \code{node} \\
  \code{node.replaceWith(...)}         & sostituisce \code{node} \\
  \code{node.remove()}                 & rimuove il nodo \\
\end{tabular}
\end{center}

\section{Interagire con l'utente}

\code{alert}, \code{prompt} e \code{confirm} sono modi semplici per interagire
con l'utente in un browser.

\begin{affianco}[0.55][c]
\begin{lstlisting}[style=html, name=interazione.html]
<ul></ul>
<script>
  alert("This page will ask for your input");

  if (confirm("Do you want to continue?") === true) {
    let name = prompt("State your name:");
    let num  = parseInt(prompt("Insert a number:"));

    for (let i = 0; i < num; i++) {
      let li = document.createElement("li");
      li.innerHTML = name;
      document.querySelector("ul").append(li);
    }
  } else {
    console.log("user cancelled");
  }
</script>
\end{lstlisting}
\risultato
\begin{immagine}[Le tre finestre, nell'ordine, e la lista\\
                 costruita con le risposte.]
\begin{tikzpicture}[
    a/.style={-{Stealth[length=4pt,width=3.5pt]}, line width=0.7pt,
      draw=neutral!45!black}]
  \node[anchor=north] (d1) at (0,0)
    {\wtdialogo{This page will ask for your input}{\wtbottone[primary!70!black][white][primary!70!black]{OK}}};
  \node[anchor=north] (d2) at (0,-2.05)
    {\wtdialogo{Do you want to continue?}{\wtbottone{Annulla}\ \wtbottone[primary!70!black][white][primary!70!black]{OK}}};
  \node[anchor=north] (d3) at (0,-4.1)
    {\wtdialogo[Ann]{State your name:}{\wtbottone{Annulla}\ \wtbottone[primary!70!black][white][primary!70!black]{OK}}};
  \draw[a] (d1.south) -- (d2.north);
  \draw[a] (d2.south) -- (d3.north);
  \node[anchor=north west, font=\Lora\small, align=left] at ([xshift=0.3cm]d1.north east)
    {\textbullet\ Ann\\ \textbullet\ Ann\\ \textbullet\ Ann};
  \node[mnota, anchor=north west] at ([xshift=0.3cm, yshift=-1.2cm]d1.north east)
    {con il numero 3};
\end{tikzpicture}
\end{immagine}
\end{affianco}

\warning{Sono finestre bloccanti}{
  \code{alert}, \code{prompt} e \code{confirm} \textbf{sospendono l'esecuzione}
  dello script e bloccano l'interazione con la pagina finché l'utente non
  risponde. Sono comodissime per imparare e per fare esperimenti, ma nelle
  applicazioni reali si preferiscono finestre di dialogo costruite nel DOM, che
  si possono impaginare, tradurre e rendere accessibili.}

\section{Eventi}

\dfn{Evento}{
  Un \term{evento} è un segnale che qualcosa è accaduto.}

Gli eventi possono riguardare:

\begin{itemize}
  \item il \textbf{comportamento dell'utente}: clic, pressioni di tasti,
        movimenti del mouse;
  \item i \textbf{form}: invio, acquisto o perdita del fuoco su un campo;
  \item il \textbf{documento}: documento o elementi caricati, richieste di rete
        completate.
\end{itemize}

\begin{center}
\renewcommand{\arraystretch}{1.35}
\begin{tabular}{@{}l@{\hspace{1.0cm}}>{\raggedright\arraybackslash}p{7.6cm}@{}}
  \textbf{Evento} & \textbf{Quando si verifica} \\[2pt]
  \multicolumn{2}{@{}l}{\color{primary!80!black}\textbf{Mouse}} \\
  \code{click}       & l'utente fa clic (o tocca) sull'elemento \\
  \code{contextMenu} & clic con il tasto destro sull'elemento \\
  \code{mouseover} / \code{mouseout} & il cursore entra o esce dall'elemento \\
  \code{mousedown} / \code{mouseup}  & il tasto principale viene premuto o
                                       rilasciato sull'elemento \\
  \code{mousemove}   & il cursore si muove \\[4pt]
  \multicolumn{2}{@{}l}{\color{primary!80!black}\textbf{Tastiera}} \\
  \code{keydown} / \code{keyup} & un tasto viene premuto o rilasciato \\[4pt]
  \multicolumn{2}{@{}l}{\color{primary!80!black}\textbf{Form}} \\
  \code{submit}      & l'utente invia un form \\
  \code{focus}       & l'utente dà il fuoco a un elemento (per esempio un
                       \tg{input}) \\[4pt]
  \multicolumn{2}{@{}l}{\color{primary!80!black}\textbf{Documento e finestra}} \\
  \code{DOMContentLoaded} & l'HTML è stato analizzato, il DOM è completo \\
  \code{load}        & la pagina è stata caricata e disegnata per intero \\
\end{tabular}
\end{center}

\subsection{Gestori di eventi}

Per reagire agli eventi si definiscono funzioni \term{gestore}
(\textit{handler}), eseguite quando l'evento scatta.

Il modo più semplice è usare gli attributi HTML \code{on<evento>}, il cui valore
è il codice JavaScript da eseguire.

\begin{affianco}[0.6][c]
\begin{lstlisting}[style=html, name=handler-attributo.html]
<input type="button" value="Click me" onclick="handleClick();">
<script>
  function handleClick() {
    console.log("Click handled");
  }
</script>
\end{lstlisting}
\risultato
\begin{browser}[handler.html]
  \wtbottone{Click me}
\end{browser}
\begin{console}[Console, dopo un clic]
Click handled
\end{console}
\end{affianco}

I gestori si possono anche definire dinamicamente da JavaScript, scrivendo la
proprietà \code{on<evento>} sull'elemento oppure usando il metodo
\code{addEventListener}.

\begin{affianco}[0.58][c]
\begin{lstlisting}[style=html, name=handler-dinamico.html]
<input type="button" value="Click me" onclick="handleClick();">
<script>
  function handleClick() { console.log("Click handled"); }

  let input = document.querySelector("input");

  input.onclick = () => { console.log("Tap"); };
  // sovrascrive il gestore definito nell'attributo HTML

  input.addEventListener("click", handleClick);
  // AGGIUNGE un nuovo gestore, senza rimuovere il precedente
</script>
\end{lstlisting}
\risultato
\begin{immagine}[\code{onclick} contiene un solo gestore;\\
                 i listener si accumulano.]
\begin{tikzpicture}
  \mchiave{1.55cm}
  \oggetto{i}{(0,0)}{3.5cm}{input}
    {onclick/() => \{ \ldots Tap \}, listener/[handleClick]}
  \node[mcod, text=neutral!50!black, anchor=west] (old) at ([xshift=6pt]i-1.east)
    {\sout{handleClick();}};
  \node[mnota, anchor=north west, text=merr] at (old.south west) {sovrascritto};
\end{tikzpicture}
\end{immagine}
\begin{console}[Console, dopo un clic]
Tap
Click handled
\end{console}
\end{affianco}

\info{Perché addEventListener è preferibile}{
  La proprietà \code{on<evento>} può contenere \textbf{un solo} gestore:
  assegnarne un altro cancella il precedente.
  \code{addEventListener} permette invece di registrare quanti gestori si
  vogliono per lo stesso evento, e di rimuoverli selettivamente con
  \code{removeEventListener}. In un'applicazione reale, dove parti diverse del
  codice possono voler reagire allo stesso evento, è l'unica scelta sensata.}

\subsection{Il valore di this in un gestore}

Dentro un gestore, \code{this} vale l'elemento a cui il gestore è stato
assegnato.

\begin{affianco}[0.6][c]
\codicepiccolo
\begin{lstlisting}[style=html, name=this-handler.html]
<input type="button" value="Click me">
<p>Hello Events!</p>
<script>
  function handleClick() {
    console.log(this.outerHTML);
  }

  document.querySelector("input").addEventListener("click", handleClick);
  // <input type="button" value="Click me">

  document.querySelector("p").addEventListener("click", handleClick);
  // <p>Hello Events!</p>
</script>
\end{lstlisting}
\risultato
\begin{browser}[this.html]
  \wtbottone{Click me}\\[6pt]
  Hello Events!
\end{browser}
\begin{console}[Console, clic sul pulsante e poi sul testo]
<input type="button" value="Click me">
<p>Hello Events!</p>
\end{console}
\end{affianco}

\subsection{L'oggetto evento}

Per gestire bene un evento servono spesso dettagli: in un \code{keydown}, quale
tasto è stato premuto? In un \code{click}, a quali coordinate? L'utente teneva
premuto \code{CTRL}?

Tutti i dettagli sono forniti in un \term{oggetto evento}, passato
automaticamente ai gestori. Tipi di evento diversi hanno proprietà diverse: per
esempio gli eventi da tastiera hanno una proprietà \code{key} che rappresenta il
tasto premuto.

\begin{affianco}[0.62][c]
\begin{lstlisting}[style=html, name=oggetto-evento.html]
<input type="text" placeholder="Write here" autofocus>
<script>
  function handler(event) {
    console.log(`You pressed ${event.key}`);
  }

  document.querySelector("input").addEventListener("keyup", handler);
</script>
\end{lstlisting}
\risultato
\begin{browser}[evento.html]
  \wtcampo[4cm][white]{JS\textbar}
\end{browser}
\begin{console}[Console, digitando J e poi S]
You pressed J
You pressed S
\end{console}
\end{affianco}

\section{Propagazione degli eventi}

\subsection{Bubbling}

\dfn{Event bubbling}{
  Quando un evento scatta su un elemento, vengono eseguiti prima tutti i gestori
  registrati su quell'elemento, poi tutti quelli registrati sul suo genitore, e
  così via risalendo fino alla radice del DOM. Questo meccanismo si chiama
  \term{event bubbling} (\guillemotleft ribollimento\guillemotright).}

Si applica alla maggior parte degli eventi, ma non a tutti: gli eventi di
\code{focus}, per esempio, non fanno bubbling.

\begin{esempio}{Il percorso di un clic}
Ogni elemento del documento ha un gestore. Facendo clic sul \tg{p}, i gestori
vengono eseguiti risalendo l'albero (a destra, nell'ordine numerato), e la
console mostra le quattro righe:

\begin{affianco}[0.55][c]
\begin{lstlisting}[style=html, name=bubbling.html]
<body>
  <section onclick="console.log('section');">
    <article onclick="console.log('article');">
      <p onclick="console.log('p');">Bubbling</p>
    </article>
  </section>
  <script>
    document.onclick = () => console.log("document");
  </script>
</body>
\end{lstlisting}
\risultato
\begin{immagine}
\begin{tikzpicture}[
    box/.style={draw, rounded corners=3pt, line width=0.7pt, font=\FiraCodeNL\scriptsize,
      anchor=north west, inner sep=0pt},
    num/.style={circle, fill=accent!70!black, text=white,
      font=\Lato\tiny\bfseries, inner sep=1pt, minimum size=10pt},
    up/.style={-{Stealth[length=5pt,width=4pt]}, line width=0.9pt, draw=accent!70!black}]
  \draw[box, draw=neutral!50!black, fill=neutral!6!white] (0,0) rectangle (5.0,-3.1);
  \node[anchor=north west, mnota] at (0.05,-0.02) {document};
  \draw[box, draw=primary!60!black, fill=primary!6!white] (0.3,-0.45) rectangle (4.7,-2.85);
  \node[anchor=north west, mcod] at (0.35,-0.47) {<section>};
  \draw[box, draw=primary!60!black, fill=primary!12!white] (0.6,-0.9) rectangle (4.4,-2.6);
  \node[anchor=north west, mcod] at (0.65,-0.92) {<article>};
  \draw[box, draw=accent!60!black, fill=accent!14!white] (0.9,-1.35) rectangle (4.1,-2.35);
  \node[anchor=north west, mcod] at (0.95,-1.37) {<p>};
  \node[font=\Lora\small] at (2.5,-1.95) {Bubbling};
  \draw[up] (3.75,-2.2) -- (3.75,0.15);
  \node[num] at (3.75,-1.6) {1};
  \node[num] at (3.75,-1.12) {2};
  \node[num] at (3.75,-0.67) {3};
  \node[num] at (3.75,-0.22) {4};
\end{tikzpicture}
\end{immagine}
\begin{console}
p
article
section
document
\end{console}
\end{affianco}

Un solo clic, quattro gestori eseguiti, dal più interno al più esterno.
\end{esempio}

\subsection{Delegazione degli eventi}

La \term{delegazione degli eventi} è un pattern potente: se molti elementi
possono generare eventi simili, invece di assegnare un gestore a ciascuno se ne
usa \textbf{uno solo} sul loro antenato comune. Gli elementi
\guillemotleft delegano\guillemotright\ la gestione all'antenato.

La proprietà \code{target} dell'oggetto evento permette di sapere su quale
elemento l'evento è stato effettivamente generato.

\begin{esempio}{Un solo gestore per tutta la tabella}
\begin{affianco}[0.58][c]
\begin{lstlisting}[style=html, name=delegazione.html]
<table onclick="handler(event);">
  <tr><th>Col</th><th>Col</th></tr>
  <tr><td>1</td><td>2</td></tr>
  <tr><td>3</td><td>4</td></tr>
  <tr><td>5</td><td>6</td></tr>
</table>
<script>
  function handler(event) {
    if (event.target.nodeName === "TD") {
      let oldValue = parseInt(event.target.innerHTML);
      event.target.innerHTML = ++oldValue;
    }
  }
</script>
\end{lstlisting}
\risultato
\begin{browser}[delegazione.html]
  \renewcommand{\arraystretch}{1.25}%
  \begin{tabular}{|c|c|}
    \hline \textbf{Col} & \textbf{Col} \\
    \hline 1 & \colorbox{accent!18!white}{3} \\
    \hline 3 & 4 \\
    \hline 5 & 6 \\ \hline
  \end{tabular}
\end{browser}
\wtnota{Dopo un clic sul 2, la cella vale 3.\\
        \code{event.target} è la \code{<td>},\\ il gestore è sulla \code{<table>}.}
\end{affianco}

Un solo gestore serve sei celle. Il vantaggio non è solo la brevità: se in
seguito si aggiungessero righe alla tabella \textbf{dinamicamente}, queste
funzionerebbero subito, senza bisogno di registrare nuovi gestori.
\end{esempio}

\subsection{Fermare il bubbling}

Un gestore può decidere che l'evento è stato completamente elaborato e che non
serve propagarlo oltre, invocando il metodo \code{stopPropagation()}
sull'oggetto evento.

\begin{affianco}[0.55][c]
\begin{lstlisting}[style=html, name=stop-propagation.html]
<body>
  <section onclick="console.log('section');">
    <article>
      <p onclick="console.log('p');">Bubbling</p>
    </article>
  </section>
  <script>
    let art = document.querySelector("article");
    art.addEventListener("click", handler);

    function handler(event) {
      console.log(this.tagName);
      event.stopPropagation();   // non arriva a <section>
    }
  </script>
</body>
\end{lstlisting}
\risultato
\begin{immagine}
\begin{tikzpicture}[
    box/.style={draw, rounded corners=3pt, line width=0.7pt, inner sep=0pt},
    up/.style={-{Stealth[length=5pt,width=4pt]}, line width=0.9pt, draw=accent!70!black}]
  \draw[box, draw=primary!60!black, fill=primary!6!white] (0.3,-0.45) rectangle (4.7,-2.85);
  \node[anchor=north west, mcod] at (0.35,-0.47) {<section>};
  \draw[box, draw=primary!60!black, fill=primary!12!white] (0.6,-0.9) rectangle (4.4,-2.6);
  \node[anchor=north west, mcod] at (0.65,-0.92) {<article>};
  \draw[box, draw=accent!60!black, fill=accent!14!white] (0.9,-1.35) rectangle (4.1,-2.35);
  \node[anchor=north west, mcod] at (0.95,-1.37) {<p>};
  \node[font=\Lora\small] at (2.5,-1.95) {Bubbling};
  \draw[up] (3.75,-2.2) -- (3.75,-0.75);
  \node at (3.75,-0.62) {\merr};
  \node[mnota, anchor=west] at (4.8,-0.62) {\code{stopPropagation()}};
\end{tikzpicture}
\end{immagine}
\begin{console}[Console, dopo un clic sul paragrafo]
p
ARTICLE
\end{console}
\end{affianco}

\warning{Fermare la propagazione con parsimonia}{
  La regola pratica è di \textbf{non} fermare la propagazione a meno di avere
  una buona ragione. Un antenato potrebbe aver bisogno di intercettare
  quell'evento, oggi o in futuro: sistemi di analisi, chiusura automatica dei
  menu a comparsa e molte librerie si basano proprio sul fatto che gli eventi
  risalgano fino in cima.}

\subsection{Generare eventi da JavaScript}

Da JavaScript è anche possibile \textbf{generare} eventi. La proprietà
\code{isTrusted} vale \code{false} per gli eventi artificiali.

\begin{affianco}[0.6][c]
\begin{lstlisting}[style=html, name=dispatch.html]
<input type="button" value="Click me">
<script>
  function handler(event) {
    console.log(`${event.type} event. Trusted: ${event.isTrusted}`);
  }

  let input = document.querySelector("input");
  input.onclick = handler;

  let click = new Event("click");
  input.dispatchEvent(click);   // isTrusted vale false
</script>
\end{lstlisting}
\risultato
\begin{console}[Console, al caricamento]
click event. Trusted: false
\end{console}
\wtnota{Un clic vero darebbe invece\\ \code{Trusted: true}.}
\end{affianco}

\section{Ordine di esecuzione degli script}

Man mano che i nostri script diventano più complessi, interagiscono con il DOM e
registrano gestori, dobbiamo sapere \textbf{quando} il codice viene eseguito
durante il caricamento di una pagina.

\begin{itemize}
  \item In generale, gli script vengono scaricati ed eseguiti \textbf{nell'ordine
        in cui compaiono}, mentre l'analisi della pagina è \textbf{sospesa}.
  \item Gli script esterni con l'attributo \attr{defer} vengono scaricati
        \textbf{in parallelo} ed eseguiti \textbf{dopo} che la pagina è stata
        completamente analizzata.
  \item Gli script esterni con l'attributo \attr{async} vengono scaricati in
        parallelo ed eseguiti \textbf{appena sono pronti}.
  \item Gli script di tipo modulo si comportano come \attr{defer} per
        impostazione predefinita, e possono essere dichiarati \attr{async}.
\end{itemize}

\begin{figure}[H]
  \centering
  \begin{tikzpicture}[font=\Lato\scriptsize,
      bar/.style={minimum height=0.4cm, anchor=west, inner sep=0pt},
      sp/.style={bar, fill=primary!22!white, draw=primary!55!black, line width=0.5pt},
      sf/.style={bar, fill=yellow-gradient-3!60!white, draw=yellow-gradient-6,
                 line width=0.5pt},
      sx/.style={bar, fill=accent!25!white, draw=accent!60!black, line width=0.5pt},
      lb/.style={font=\FiraCodeNL\scriptsize, anchor=east,
                 text=neutral!25!black}]

    % ─── legenda ───
    \node[sp, minimum width=0.7cm] at (0,1.35) {};
    \node[font=\Lato\scriptsize, anchor=west, text=neutral!30!black] at (0.8,1.35)
      {analisi della pagina};
    \node[sf, minimum width=0.7cm] at (3.5,1.35) {};
    \node[font=\Lato\scriptsize, anchor=west, text=neutral!30!black] at (4.3,1.35)
      {download dello script};
    \node[sx, minimum width=0.7cm] at (7.3,1.35) {};
    \node[font=\Lato\scriptsize, anchor=west, text=neutral!30!black] at (8.1,1.35)
      {esecuzione};

    % ─── momento in cui il tag viene incontrato ───
    \draw[dashed, draw=neutral!45!white, line width=0.6pt] (1.2,0.65) -- (1.2,-5.0);
    \node[font=\Lato\scriptsize, text=neutral!35!black, anchor=south] at (1.2,0.7)
      {il tag \code{script} viene incontrato};

    % ─── 1: script classico ───
    \node[lb] at (-0.2,0.2) {<script>};
    \node[sp, minimum width=1.2cm] at (0,0.2) {};
    \node[sf, minimum width=1.6cm] at (1.2,0.2) {};
    \node[sx, minimum width=1.0cm] at (2.8,0.2) {};
    \node[sp, minimum width=2.2cm] at (3.8,0.2) {};

    % ─── 2: defer ───
    \node[lb] at (-0.2,-0.9) {defer};
    \node[sp, minimum width=1.2cm] at (0,-0.9) {};
    \node[sp, minimum width=3.4cm] at (1.2,-0.9) {};
    \node[sf, minimum width=1.6cm] at (1.2,-1.35) {};
    \node[sx, minimum width=1.0cm] at (4.6,-0.9) {};

    % ─── 3: async ───
    \node[lb] at (-0.2,-2.0) {async};
    \node[sp, minimum width=1.2cm] at (0,-2.0) {};
    \node[sp, minimum width=1.6cm] at (1.2,-2.0) {};
    \node[sf, minimum width=1.6cm] at (1.2,-2.45) {};
    \node[sx, minimum width=1.0cm] at (2.8,-2.0) {};
    \node[sp, minimum width=2.2cm] at (3.8,-2.0) {};

    % ─── 4: script di tipo modulo ───
    \node[lb] at (-0.2,-3.1) {type="module"};
    \node[sp, minimum width=1.2cm] at (0,-3.1) {};
    \node[sp, minimum width=3.4cm] at (1.2,-3.1) {};
    \node[sf, minimum width=1.6cm] at (1.2,-3.55) {};
    \node[sx, minimum width=1.0cm] at (4.6,-3.1) {};

    % ─── 5: modulo async ───
    \node[lb] at (-0.2,-4.2) {type="module" async};
    \node[sp, minimum width=1.2cm] at (0,-4.2) {};
    \node[sp, minimum width=1.6cm] at (1.2,-4.2) {};
    \node[sf, minimum width=1.6cm] at (1.2,-4.65) {};
    \node[sx, minimum width=1.0cm] at (2.8,-4.2) {};
    \node[sp, minimum width=2.2cm] at (3.8,-4.2) {};
  \end{tikzpicture}
  \caption{Quando viene scaricato ed eseguito uno script, a seconda degli
           attributi.}
  \label{fig:script-order}
\end{figure}

\begin{esempio}{Perché l'ordine conta}
Consideriamo lo script \code{load.js} (a sinistra), che cerca un \tg{h1} nel
\tg{body}, e la pagina che lo include nell'\tg{head} (a destra):

\begin{affianco}[0.45]
\begin{lstlisting}[style=js, name=load.js]
let h1 = document.body.querySelector("h1");
console.log(h1.innerText);
\end{lstlisting}
\risultato
\begin{lstlisting}[style=html, name=pagina.html]
<!DOCTYPE html>
<html>
  <head>
    <meta charset="utf-8">
    <title>Events</title>
    <script src="load.js"></script>
  </head>
  <body>
    <h1>JavaScript</h1>
  </body>
</html>
\end{lstlisting}
\end{affianco}

\begin{center}
\renewcommand{\arraystretch}{1.4}
\begin{tabular}{@{}l@{\hspace{1.1cm}}>{\raggedright\arraybackslash}p{8.6cm}@{}}
  \textbf{Variante} & \textbf{Esito} \\
  \code{<script src="load.js">}
    & \textbf{Errore.} Lo script viene eseguito subito, quando il \tg{body} non
      è ancora stato analizzato: \code{h1} vale \code{null}. \\
  \code{<script defer src="load.js">}
    & \textbf{Funziona.} L'esecuzione avviene dopo l'analisi completa della
      pagina. \\
  \code{<script async src="load.js">}
    & \textbf{Non si può sapere.} Dipende da quanto tempo serve a scaricare lo
      script. \\
\end{tabular}
\end{center}

L'alternativa che funziona sempre, indipendentemente da dove si trova il tag, è
attendere l'evento \code{DOMContentLoaded}:

\begin{lstlisting}[style=js, name=load-sicuro.js]
document.addEventListener("DOMContentLoaded", () => {
  let h1 = document.body.querySelector("h1");
  console.log(h1.innerText);
});
\end{lstlisting}
\end{esempio}

\section{Riferimenti}

\begin{itemize}
  \item \textbf{The Modern JavaScript Tutorial} \\
        \urlx{https://javascript.info/} \\
        Parte 2: \textit{Document} (1.1--1.7), \textit{Introduction to events},
        \textit{UI events}, \textit{Document and resource loading}.
  \item \textbf{Eloquent JavaScript} (3\textsuperscript{a} edizione), Marijn
        Haverbeke. \\
        \urlx{https://eloquentjavascript.net/} \\
        Capitoli da 13 a 15.
  \item \textbf{JavaScript Reference} --- MDN web docs. \\
        \urlx{https://developer.mozilla.org/en-US/docs/Web/JavaScript/Reference}
\end{itemize}
